\section{Problem Restatement}

The problem asks the paper to recommend an operation scheme for an urban rail corridor and to show that the resulting timetable is still operationally credible. This means the task is not satisfied by a passenger-flow description alone, nor by a timetable drawing alone. The paper must connect demand concentration, service-pattern choice, and timetable feasibility in one continuous chain.

For a contest-style draft, the task is naturally decomposed into three layers. The first layer identifies where passenger demand is concentrated and whether the overlap zone is strong enough to justify a large-small-route service structure. The second layer determines the service pattern itself, including route division, headway, and train allocation. The third layer converts the selected operation plan into a station-by-station timetable and checks whether the outcome respects tracking, dwell, and capacity constraints. This decomposition should be visible early, because it decides the page rhythm and evidence order of the entire paper.

Accordingly, the core deliverables of the paper must include more than one planning table. A complete draft should show at least a section-flow or OD-based evidence artifact, an operation-plan comparison artifact, a selected-plan table, a machine-readable timetable output, and one or more feasibility-audit artifacts. Only with this set of outputs can the final recommendation be read as a route from demand to service to schedule rather than as a one-step optimization claim.

The restatement is also important for later writing style. Once the problem is defined as a chain from flow analysis to service design to timetable audit, the paper can maintain a stable question map from abstract to conclusion. That stability is one of the easiest ways to make the draft feel like a serious human-team paper instead of a stack of disconnected technical blocks.
